% =========================================================
% Rational Density and Order Facts
% =========================================================

\subsection{Rational Density and Order Facts}
\label{sec:rational-density-and-order}

The rational numbers are not merely numerous; they are tightly packed. Between
any two distinct rationals one can find another, and small rational
perturbations can always be made without leaving $\mathbb{Q}$.

% ---------------------------------------------------------
% Theorem: Between Any Two Rational Numbers Lies Another Rational
% ---------------------------------------------------------
\begin{tcolorbox}[colback=thmbox, colframe=thmborder, arc=2pt,
  left=6pt, right=6pt, top=4pt, bottom=4pt,
  title={\small\textbf{Theorem (Between Any Two Rational Numbers Lies Another Rational)}},
  fonttitle=\small\bfseries]
\begin{theorem}[Between Any Two Rational Numbers Lies Another Rational]
\label{thm:between-any-two-rationals-is-a-rational}
If $r,s\in\mathbb{Q}$ and $r<s$, then there exists $t\in\mathbb{Q}$ such that
\[
r<t<s.
\]
\hyperref[prf:between-any-two-rationals-is-a-rational]{Go to proof.}
\end{theorem}
\end{tcolorbox}

\begin{remark*}[Standard quantified statement]
\[
\forall r,s\in\mathbb{Q}\;
\bigl(r<s \Longrightarrow \exists t\in\mathbb{Q}\;(r<t<s)\bigr).
\]
\end{remark*}

\begin{remark*}[Predicate reading]
\[
\operatorname{Rational}(r)\land \operatorname{Rational}(s)\land r<s
\Longrightarrow
\exists t\;\bigl(\operatorname{Rational}(t)\land r<t<s\bigr).
\]
\end{remark*}

\begin{remark*}[Negated quantified statement]
\[
\exists r,s\in\mathbb{Q}\;
\bigl(r<s \land \forall t\in\mathbb{Q}\;(t\le r \lor s\le t)\bigr).
\]
\end{remark*}

\begin{remark*}[Negation predicate reading]
\[
\operatorname{Rational}(r)\land \operatorname{Rational}(s)\land r<s
\land
\neg\exists t\;\bigl(\operatorname{Rational}(t)\land r<t<s\bigr).
\]
\end{remark*}

\begin{remark*}[Failure modes]
The theorem could fail only if some pair of distinct rationals had no room for
another rational between them. The midpoint construction shows this never
happens in $\mathbb{Q}$.
\end{remark*}

\begin{remark*}[Failure mode decomposition]
\[
r<s
\Longrightarrow
r<\frac{r+s}{2}<s,
\qquad
\frac{r+s}{2}\in\mathbb{Q}.
\]
\end{remark*}

\begin{remark*}[Interpretation]
This is the order-density of $\mathbb{Q}$ in itself. Rational numbers have no
adjacent neighbors.
\end{remark*}

\begin{remark*}[Dependencies]
\begin{itemize}
  \item \hyperref[thm:q-is-ordered-field]{$\mathbb{Q}$ Is an Ordered Field}
\end{itemize}
\end{remark*}

% ---------------------------------------------------------
% Corollary: Q Is Dense in Itself
% ---------------------------------------------------------
\begin{corollary}[$\mathbb{Q}$ Is Dense in Itself]
\label{cor:q-dense-in-itself}
The subset $\mathbb{Q}$ is order dense in itself.
\hyperref[prf:q-dense-in-itself]{Go to proof.}
\end{corollary}

\begin{remark*}[Standard quantified statement]
\[
\forall r,s\in\mathbb{Q}\;
\bigl(r<s \Longrightarrow \exists t\in\mathbb{Q}\;(r<t<s)\bigr).
\]
\end{remark*}

\begin{remark*}[Predicate reading]
\[
\operatorname{DenseIn}(\mathbb{Q},\mathbb{Q}).
\]
\end{remark*}

\begin{remark*}[Negated quantified statement]
\[
\exists r,s\in\mathbb{Q}\;
\bigl(r<s \land \forall t\in\mathbb{Q}\;(t\le r \lor s\le t)\bigr).
\]
\end{remark*}

\begin{remark*}[Negation predicate reading]
\[
\neg\operatorname{DenseIn}(\mathbb{Q},\mathbb{Q}).
\]
\end{remark*}

\begin{remark*}[Failure modes]
This could fail only if $\mathbb{Q}$ contained a genuine gap when viewed as an
ordered set in its own right. The previous theorem rules that out.
\end{remark*}

\begin{remark*}[Dependencies]
\begin{itemize}
  \item \hyperref[thm:between-any-two-rationals-is-a-rational]{Between Any Two Rational Numbers Lies Another Rational}
\end{itemize}
\end{remark*}

% ---------------------------------------------------------
% Corollary: The Rational Numbers Have No Adjacent Points
% ---------------------------------------------------------
\begin{corollary}[The Rational Numbers Have No Adjacent Points]
\label{cor:q-has-no-adjacent-points}
There do not exist rational numbers $r<s$ such that no rational number lies
strictly between them.
\hyperref[prf:q-has-no-adjacent-points]{Go to proof.}
\end{corollary}

\begin{remark*}[Standard quantified statement]
\[
\neg\exists r,s\in\mathbb{Q}\;
\bigl(r<s \land \forall t\in\mathbb{Q}\;(t\le r \lor s\le t)\bigr).
\]
\end{remark*}

\begin{remark*}[Predicate reading]
\[
\forall r,s\in\mathbb{Q}\;
\bigl(r<s \Longrightarrow \exists t\in\mathbb{Q}\;(r<t<s)\bigr).
\]
\end{remark*}

\begin{remark*}[Negated quantified statement]
\[
\exists r,s\in\mathbb{Q}\;
\bigl(r<s \land \neg\exists t\in\mathbb{Q}\;(r<t<s)\bigr).
\]
\end{remark*}

\begin{remark*}[Failure modes]
Adjacent points would mean that some rational interval had no internal rational
point. The midpoint argument prevents that.
\end{remark*}

\begin{remark*}[Dependencies]
\begin{itemize}
  \item \hyperref[thm:between-any-two-rationals-is-a-rational]{Between Any Two Rational Numbers Lies Another Rational}
\end{itemize}
\end{remark*}

% ---------------------------------------------------------
% Theorem: Q Has No Minimum Positive Element
% ---------------------------------------------------------
\begin{tcolorbox}[colback=thmbox, colframe=thmborder, arc=2pt,
  left=6pt, right=6pt, top=4pt, bottom=4pt,
  title={\small\textbf{Theorem ($\mathbb{Q}$ Has No Minimum Positive Element)}},
  fonttitle=\small\bfseries]
\begin{theorem}[$\mathbb{Q}$ Has No Minimum Positive Element]
\label{thm:q-has-no-minimum-positive-element}
For every $q\in\mathbb{Q}$ with $q>0$, there exists $r\in\mathbb{Q}$ such that
\[
0<r<q.
\]
\hyperref[prf:q-has-no-minimum-positive-element]{Go to proof.}
\end{theorem}
\end{tcolorbox}

\begin{remark*}[Standard quantified statement]
\[
\forall q\in\mathbb{Q}\;
\bigl(q>0 \Longrightarrow \exists r\in\mathbb{Q}\;(0<r<q)\bigr).
\]
\end{remark*}

\begin{remark*}[Predicate reading]
\[
\operatorname{Rational}(q)\land q>0
\Longrightarrow
\exists r\;\bigl(\operatorname{Rational}(r)\land 0<r<q\bigr).
\]
\end{remark*}

\begin{remark*}[Negated quantified statement]
\[
\exists q\in\mathbb{Q}\;
\bigl(q>0 \land \forall r\in\mathbb{Q}\;(0<r \Longrightarrow q\le r)\bigr).
\]
\end{remark*}

\begin{remark*}[Failure modes]
This would fail only if some positive rational were the smallest positive
rational. Halving it produces a strictly smaller positive rational, so no such
minimum can exist.
\end{remark*}

\begin{remark*}[Interpretation]
Positive rationals can always be halved. There is no smallest positive
rational number.
\end{remark*}

\begin{remark*}[Dependencies]
\begin{itemize}
  \item \hyperref[thm:q-is-ordered-field]{$\mathbb{Q}$ Is an Ordered Field}
\end{itemize}
\end{remark*}

% ---------------------------------------------------------
% Lemma: Upper/Lower Perturbation Lemma (Rational Version)
% ---------------------------------------------------------
\begin{lemma}[Upper/Lower Perturbation Lemma, Rational Version]
\label{lem:rational-upper-lower-perturbation}
If $q\in\mathbb{Q}$ and $\varepsilon\in\mathbb{Q}$ with $\varepsilon>0$,
then there exist $r,s\in\mathbb{Q}$ such that
\[
q-\varepsilon < r < q < s < q+\varepsilon.
\]
\hyperref[prf:rational-upper-lower-perturbation]{Go to proof.}
\end{lemma}

\begin{remark*}[Standard quantified statement]
\[
\forall q,\varepsilon\in\mathbb{Q},\;\varepsilon>0\;\;
\exists r,s\in\mathbb{Q}\;
\bigl(q-\varepsilon < r < q < s < q+\varepsilon\bigr).
\]
\end{remark*}

\begin{remark*}[Negated quantified statement]
\[
\exists q,\varepsilon\in\mathbb{Q},\;\varepsilon>0\;\;
\forall r,s\in\mathbb{Q}\;
\neg\bigl(q-\varepsilon < r < q < s < q+\varepsilon\bigr).
\]
\end{remark*}

\begin{remark*}[Interpretation]
Every rational point has rational points strictly on both sides within any
prescribed positive distance. This is the working tool for punctured-ball
arguments in $\mathbb{Q}$.
\end{remark*}

\begin{remark*}[Dependencies]
\begin{itemize}
  \item \hyperref[thm:q-has-no-minimum-positive-element]{$\mathbb{Q}$ Has No Minimum Positive Element}
\end{itemize}
\end{remark*}

% ---------------------------------------------------------
% Lemma: Rational Perturbation Lemma
% ---------------------------------------------------------
\begin{lemma}[Rational Perturbation Lemma]
\label{lem:rational-perturbation-lemma}
If $q\in\mathbb{Q}$ and $\varepsilon>0$, then there exists $r\in\mathbb{Q}$
such that
\[
0<|r-q|<\varepsilon.
\]
\hyperref[prf:rational-perturbation-lemma]{Go to proof.}
\end{lemma}

\begin{remark*}[Standard quantified statement]
\[
\forall q\in\mathbb{Q},\;\forall\varepsilon>0\;\;
\exists r\in\mathbb{Q}\;
\bigl(0<|r-q|<\varepsilon\bigr).
\]
\end{remark*}

\begin{remark*}[Negated quantified statement]
\[
\exists q\in\mathbb{Q},\;\exists\varepsilon>0\;\;
\forall r\in\mathbb{Q}\;
\bigl(r=q \lor |r-q|\ge\varepsilon\bigr).
\]
\end{remark*}

\begin{remark*}[Interpretation]
Every rational point is a limit point of $\mathbb{Q}$: it can be approached
arbitrarily closely by other rational numbers.
\end{remark*}

\begin{remark*}[Dependencies]
\begin{itemize}
  \item \hyperref[lem:rational-upper-lower-perturbation]{Upper/Lower Perturbation Lemma}
\end{itemize}
\end{remark*}

% ---------------------------------------------------------
% Theorem: Rational Chain with Uniformly Small Steps
% ---------------------------------------------------------
\begin{tcolorbox}[colback=thmbox, colframe=thmborder, arc=2pt,
  left=6pt, right=6pt, top=4pt, bottom=4pt,
  title={\small\textbf{Theorem (Rational Chain with Uniformly Small Steps)}},
  fonttitle=\small\bfseries]
\begin{theorem}[Rational Chain with Uniformly Small Steps]
\label{thm:rational-chain-with-small-steps}
Let $r,s\in\mathbb{Q}$ with $r<s$, and let $p\in\mathbb{Q}$ with $p>0$.
Then there exist $m\in\mathbb{N}$ and rational numbers
$q_1,q_2,\dots,q_m$ such that
\[
q_1=r,\qquad q_m=s,
\]
and
\[
0<q_k-q_{k-1}<p
\qquad\text{for every }k=2,\dots,m.
\]
\hyperref[prf:rational-chain-with-small-steps]{Go to proof.}
\end{theorem}
\end{tcolorbox}

\begin{remark*}[Standard quantified statement]
\[
\forall r,s,p\in\mathbb{Q}\;
\Bigl(
r<s \land p>0
\Longrightarrow
\exists m\in\mathbb{N}\;\exists q_1,\dots,q_m\in\mathbb{Q}\;
\bigl(q_1=r \land q_m=s \land
\forall k\in\{2,\dots,m\}\;(0<q_k-q_{k-1}<p)\bigr)
\Bigr).
\]
\end{remark*}

\begin{remark*}[Interpretation]
Any rational interval can be traversed by finitely many rational steps each
smaller than a prescribed positive bound. This is the discrete counterpart of
the Archimedean property applied to step traversal.
\end{remark*}

\begin{remark*}[Dependencies]
\begin{itemize}
  \item \hyperref[thm:archimedean-q]{Archimedean Property of $\mathbb{Q}$}
\end{itemize}
\end{remark*}
